\subsubsection{Classical specializations of both Hecke quotients}
\label{par:count-two-classical}
\hypertarget{source-ordinary-count-ii-point-output}{}

For a characteristic-zero target the two ordinary evaluations yield representations. To compare their Frobenius polynomials with those of the original point, the evaluations must first be realized at classical weight.

\begin{proposition}[Classical control at the propagated target]\label{prop:two-link-classical-control}
At a characteristic-zero even-order target reached by the
two-link propagation, each evaluation factors through a classical
ordinary eigensystem in its link. Let \(r_i\) be the Galois representation of the eigensystem at level \(i\), \(r\) the representation at the chosen determinant point and \(\kappa\) the common twisting character. Then 
\[
 \det(X-r_i(g))=\det(X-\kappa(g)r(g))\qquad(g\in G_L).
\]

\end{proposition}

\begin{proof}
Use the same Hecke quotients \(R_i/J_i\) as in
Paragraph~\ref{par:count-two-ideals}. At a characteristic-zero
target \(y\) reached by Paragraph~\ref{par:count-two-propagation},
the evaluations \(x_i=y|_B f_i\) factor through these quotients.
Write \(m=(j+1,j,1,0)\) for the rational inertial exponents.
On the selected inertia groups above \(\ell\), \(\kappa\) is
an integral cyclotomic power up to a finite character. It
adds a common integer to the four exponents at each place.
Thus the coefficients have the algebraic parts given by
\(\lambda_{\tau,3-i}=-m_{w,i}-i\), with \(m_{w,i}\) the exponents
after twisting. The difference in the highest weight order is
\[
       \lambda_{\tau,2-i}-\lambda_{\tau,3-i}
                  =m_{w,i}-m_{w,i+1}-1\ge0 .
\]
The difference inequality gives dominance. After the algebraic
part is removed the remaining compact torus character is smooth. These conditions put the point in the arithmetic ordinary locus used in the cohomology calculation.

Here is the control argument at this weight. Write \(O_x\)
for a coefficient DVR containing the weight values, \(\pi_x\)
for its uniformizer, \(W_\lambda\) for the algebraic coefficient
representation, and \(p_\lambda\) for its lowest weight quotient
for the upper unipotent group. Use a lattice in \(W_\lambda\)
stable under the integral local groups and decomposed into
algebraic torus weights. Such a lattice is given by the integral
highest weight construction at the embeddings, tensored with
the required integral determinant powers. Write \(L\subset W_\lambda\) for this integral lattice and \(L_\mu\) for its weight submodules. Then
\[
 L=L_{\rm low}\oplus\ker(p_\lambda|_L),\qquad L_{\rm low}\simeq O_x,
 \qquad \ker(p_\lambda|_L)=\bigoplus_{\mu\ne\mu_{\rm low}}L_\mu .
\]
This gives the rank-one projection integrally.
Choose \(c\) so that the smooth character kills \(T_c\) and
\(\bar N_c\) acts trivially on the coefficient lattice modulo \(\pi_x\).
Use \(J_c=NT\bar N_c\) for the following comparison.

Trivial stabilizers identify the classical module with the
finite module of values on double coset representatives.
In the left-invariant realization \(F_\ell\) of
Section~\ref{sec:count-two-seed-specialization}, the dominant
translation \(z\) acts by the sum with coefficient matrices
\(W_\lambda(b)W_\lambda(z)\), where \(b\) ranges over
\(N/zNz^{-1}\). Normalize each matrix by the lowest-weight
scalar of \(z\). The normalized \(W_\lambda(z)\) is integral
and kills the kernel of \(p_\lambda\) modulo \(\pi_x\):
For any other weight \(\mu\), write
\(\mu-\mu_{\rm low}=\sum_\alpha n_\alpha\alpha\) with
\(n_\alpha\ge0\), some \(n_\alpha>0\).

Then
\[
 v_{\pi_x}\!\left(\frac{\mu(z)}{\mu_{\rm low}(z)}\right)
      =\sum_\alpha n_\alpha v_{\pi_x}(\alpha(z))>0.
\]
Modulo \(\pi_x\), the normalized diagonal operator kills all
these higher weight spaces. The integral unipotent matrices
\(W_\lambda(b)\) preserve the quotient projection, so the same
holds for the normalized translation sum. The resulting operator intertwines the
ordinary operator on the scalar functions under \(p_\lambda\).
On characteristic-zero functions the projection is into the
scalar module with lower unipotent depth allowed to increase
with precision, by continuity of the coefficient action.

Modulo \(\pi_x\), the scalar module at depth \(c\) is obtained
by projecting the algebraic values on the same cosets. Let \(P:C\to S\) be this projection on the
depth-\(c\) residue spaces. It is onto, commutes with the
normalized translation \(U\), and \(U\ker P=0\).
Since \(P\) is equivariant and onto, \(P(U^a C)=U^a S\)
for every \(a\).

At the stable images its kernel is
\[
    C^{\rm ord}\cap\ker P \ \subseteq\ C^{\rm ord}\cap\ker U=0 .
\]
The last equality uses invertibility on \(C^{\rm ord}\).
This proves the comparison on the stable images. The integral ordinary projectors reduce
to these stable-image projectors. The scalar ordinary module
is finite free with the value filtration of
Paragraph~\ref{par:ordinary-functions}; its ordinary reduction
is exactly the stable image at depth \(c\). Thus the projected
classical ordinary module maps isomorphically onto the scalar
ordinary module, by freeness and reduction modulo \(\pi_x\).
The spherical Hecke actions and their residual localizations
commute with the projection. Together with the scalar weight
base change of Section~\ref{sec:count-two-seed-specialization},
this gives the classical realization on each \(M_i\).

Any field evaluation of \(R_i/J_i\) at this weight occurs in
that specialization. Indeed, putting \(A_i=R_i/J_i\), finiteness and faithfulness imply
\[
 \operatorname{Supp}_{A_i}M_i=V(\operatorname{Ann}_{A_i}M_i)
 =\operatorname{Spec}A_i,\qquad M_i\otimes_{A_i}\kappa(\mathfrak p)\ne0 .
\]
The last assertion is Nakayama's lemma at the evaluation prime \(\mathfrak p\).
After extending coefficients, that fiber has the evaluation
as an eigensystem. In the automorphic decomposition choose a
contributing constituent with these spherical eigenvalues.
Its semisimple Galois attachment has the characteristic
polynomials of \(\kappa\,r|_{G_L}\), by the spherical
normalization on the dense set of split good Frobenius classes.
This equality holds on all of \(G_L\), including at reducible
targets. At an irreducible target it identifies the actual
Galois representations.
\hfill\mbox{\cite{newton-thorne-unitary}}\label{cor:ordinary-count-ii-point-output-1}

\end{proof}

\subsubsection{The Frobenius ratio test at a full block}
\label{par:count-two-purity}

A hypothetical full two-dimensional block imposes a determinant ratio at good primes. The polynomials just obtained from both ordinary levels give an incompatible restriction on their roots.

\begin{proposition}[Exclusion of a full block]\label{prop:count-two-full-block-excluded}
At a characteristic-zero full-block target, write \(\alpha_i(w)\) for the Frobenius roots at a good prime in the block test. The two-link specialization gives 
\[
 |\alpha_i(w)/\alpha_j(w)|<Nw\qquad\text{for every }i,j.
\]

This contradicts the determinant ratio of the two blocks.
\end{proposition}

\begin{proof}
In the transfer of a contributing constituent from link \(i\),
a discrete Speh summand of length at least two would give two
summands of its Galois representation differing by
\(\epsilon_\ell\) or its inverse. Take the chosen Frobenius
element \(g_i\) at the exceptional place of this link.
All four roots of the characteristic polynomial at \(g_i\)
reduce to the same nonzero residual eigenvalue, by the
construction of \(J_i\) and \(f_i\). The eigenvalues of both
summands are units, since they are roots of that integral
polynomial with unit constant coefficient. If \(z\) and \(z\epsilon_\ell(g_i)^{\pm1}\) are such roots,
then \(z\) is a unit. Their reductions cannot agree:
\[
              \bar z\,\bar\epsilon_\ell(g_i)^{\pm1}\ne\bar z .
\]
This contradicts the common nonzero residual eigenvalue.

Consequently every discrete summand is cuspidal. The
cuspidal summands in the unitary base change are tempered
at the good places, by the results cited after
Corollary~1.3 of Newton and Thorne. The common algebraic
normalization multiplies every root at a place by the same
power of its residue norm. Thus all roots have the same
absolute value there, and in particular each ordered ratio
of absolute values is strictly less than the residue norm.

The purity condition for that calculation is the ratio condition just proved. It passes to the rational representation:
If \(\alpha_i\) are the rational roots, restriction to a
place \(w\mid p\) of residue degree \(f\) gives \(\alpha_i^f\).
The common scalar twist cancels. The inequalities over \(L\) imply
\[
              |\alpha_i/\alpha_j|^f<Nw=p^f
            \quad\Longrightarrow\quad|\alpha_i/\alpha_j|<p.
\] The property
also holds at primes splitting completely in \(L\), which
suffices for the block-determinant test. Apply the link to
the nonsplit representative constructed in
Paragraph~\ref{par:full-block-extension}. Equality of polynomials
gives the same ratio bound at the split block point.
\hfill\mbox{\cite{newton-thorne-unitary}}\label{cor:ordinary-count-ii-point-output-2}

\end{proof}

\subsubsection{Attachment at an irreducible target}
\label{par:count-two-glue}

At an irreducible target, equality of the characteristic polynomials identifies the restricted representations themselves. Using both levels then compares their local parameters outside their respective exceptional sets.

\begin{proposition}[Cuspidal attachment from the two links]\label{prop:count-two-cuspidal-attachment}
At a rational irreducible target of the even-order construction,
the restrictions to \(G_L\) are absolutely irreducible. The
cuspidal transfers from both ordinary links give the same
representation there; if \(S_1,S_2\) are the respective sets of exceptional places, the compatible local parameters cover both choices because 
\[
 S_1\cap S_2=\varnothing.
\]

\end{proposition}

\begin{proof}
For a rational irreducible target, restriction to \(G_L\)
is absolutely irreducible. The restriction to a normal
subgroup is semisimple, by taking translates of a simple
subspace. If the restriction were reducible, the action of
the abelian quotient \(\operatorname{Gal}(L/\Q)\) on its
commutant would have a nontrivial character. An eigenvector in the commutant is a nonzero intertwiner
for a self-twist \(\chi\), hence invertible by irreducibility on
the rational group. It yields
\[
              (1-\chi(g))\operatorname{tr}r(g)=0 .
\]
On any coset where \(\chi\ne1\) the trace is zero. The reduction of the trace on such a coset
is a linear combination of the four distinct restricted
residual labels with nonzero coefficients. Linear independence
of characters rules out its identical vanishing. This proof
also applies to any target obtained by the change of frame
or scalar extension in the modular-locus argument.

Each link now gives a cuspidal transfer over \(L\): an isobaric
sum with several terms would make the attachment reducible.
At a common good prime the same twist and normalization give
\[
     P_{\Pi_1,w}(T)=\det(T-\kappa r(\mathrm{Frob}_w))
                         =P_{\Pi_2,w}(T),
\]
where \(P_{\Pi_i,w}\) denotes the algebraically normalized local
polynomial. Both links use the same specialization \(y\). Strong multiplicity
one identifies the transfers. For each selected place, use
the local comparison at a level where it is outside the
exceptional pair of places over \(L^+\). Indeed if \(S_i\) is the exceptional pair over the \(i\)-th choice of \(v_0\), their rational primes are distinct. Consequently
\[
 S_1\cap S_2=\varnothing,\qquad S_1^c\cup S_2^c
 =\{\text{all finite places}\}.
\]

At each of the exceptional places the other link therefore supplies the local comparison. At the other places outside the level the
transfer is unramified, and at the marked places above
\(\ell\) the split local comparison gives the Weil--Deligne
parameter of the same target.
The common scalar \(\kappa\) is locally algebraic at \(\ell\)
with integral labeled exponents and finite inertia away from
\(\ell\). By reciprocity it is the \(\ell\)-adic avatar of an
algebraic Hecke character. Twisting the transfer by its inverse
gives the cuspidal attachment of \(r|_{G_L}\), with the rational
similitude as its multiplier. In particular the target weights
are recovered after removing the twist.
\hfill\mbox{\cite{newton-thorne-unitary}}\label{cor:ordinary-count-ii-point-output-3}

\end{proof}

\subsubsection{Ordinary flag and rational level}
\label{par:count-two-host-level}

The attachment is obtained after restriction to the fitted field. To count at a prescribed rational level, the original representation must have the required ordered inertia exponents and conductor bounds.

\begin{proposition}[The rational ordinary flag and its level]\label{prop:count-two-rational-level}
For an irreducible even-order target with \(j\ge2\) in
Section~\ref{sec:odd-prime-cases} and the attachment of
Proposition~\ref{prop:count-two-cuspidal-attachment}, the rational
ordinary flag has the prescribed exponents and zero monodromy.
The rational inertia conditions give the level bound for the ensuing quadratic restriction, with 
\[
 M=\ell\,2^{n_{\rm II}}\mathcal N.
\]

\end{proposition}

\begin{proof}
\proofstep{Recover the rational ordinary characters.}

At \(\ell\) write the ordered characters of \(r\) as
\(\kappa_i=\nu_i\epsilon_\ell^{m_i}\) with
\(m=(j+1,j,1,0)\) and \(\nu_i\) finite on inertia. The
character \(\kappa_i/\kappa_k\) can equal the cyclotomic
character on a full decomposition group only if
\(m_i-m_k=1\). Restrict to an open inertia subgroup
killing \(\nu_i/\nu_k\) and use the infinite cyclotomic image
there to see this. As \(j\ge2\), the consecutive gaps are
\((1,j-1,1)\), each at least one.

If \(m_i-m_k=1\) with
\(i<k\), then \(k=i+1\). At the retained Frobenius lift of
Paragraph~\ref{par:count-ii-tower}, the cyclotomic
character and the residual \(\bar\Psi\) are one, and
\(\bar\eta\) has order greater than four. In both permitted frames the labels at this lift are
\[
         (\bar\eta,1,\bar\eta,1)\qquad(\bar\eta\ne1).
\]
Thus all adjacent labels differ at a point where
\(\bar\epsilon_\ell=1\). This rules out the cyclotomic equality at the
marked place of \(L\) as well as over \(\Q_\ell\).

The target is de Rham at \(\ell\), since it is so over the
fitted CM field and de Rhamness descends along a finite
extension. Its local Weil--Deligne representation is
therefore defined. By exactness on the ordinary flag its
graded Weil lines, on the rational Weil group, are
\[
       \mu_i(s)=\iota(\nu_i(s))\,\ell^{m_i\deg(s)}
\]
in the arithmetic Frobenius convention, where \(\deg(s)\)
is the exponent of arithmetic Frobenius in the unramified
quotient.

At the marked place of \(L\) use the restriction
of these lines. A norm ratio between two lines would give
\(\kappa_i/\kappa_k=\epsilon_\ell\) in one order: taking
\(\ell\)-adic valuations at Frobenius forces the exponent
difference to be one, because the \(\nu_i\) take unit
values. Equality then forces the two \(\nu_i\) to agree
by density of the Weil group. The Weil--Deligne relation on the \((i,k)\) entry of \(N\)
forces a norm ratio \(\mu_i/\mu_k=|\cdot|^{\pm1}\) if
that entry survives. Indeed the ratio constraint follows entrywise from the Weil--Deligne identity
\[
 (\mu_i(s)/\mu_k(s)-|s|)N_{ik}=0
\]
in the convention \(r(s)Nr(s)^{-1}=|s|N\). Switching conventions inverts the norm. Since both orders have been excluded, \(N=0\) over \(\Q_\ell\) and after restriction.

At the marked place of \(L\) all \(|\mu_i/\mu_k|\) on
Weil elements are one, by temperedness of the glued
cuspidal transfer and its local compatibility. Removing
the scalar twist preserves this equality. It holds on the
rational Weil group by taking positive powers into this
local subgroup. Choose \(s\) in the rational Weil group
mapping to the uniformizer \(\ell\) under reciprocity. Then
\(\epsilon_\ell(s)=1\) and \(\deg(s)\ne0\). Were \(\kappa_i(s)/\kappa_k(s)\) torsion, its complex modulus
would be one.

Since \(\epsilon_\ell(s)=1\), the displayed
Weil lines would then give
\[
       |\mu_i(s)/\mu_k(s)|=\ell^{(m_i-m_k)\deg(s)}\ne1 .
\]
Temperedness on the restricted Weil group forces absolute value one. If \(s^a\) lies there, then
\[
 |(\mu_i/\mu_k)(s^a)|=|(\mu_i/\mu_k)(s)|^a .
\]
Thus the violation persists on restriction, a contradiction. The character ratios in question have infinite order. A power of \(s\) fixes \(L\), and
further positive powers give the same conclusions over
finite local extensions. This element will be used for
adjoint rigidity.

Away from \(\ell\) the rational chart kills an open subgroup
of inertia. At primes of \(S_b\) it imposes the inertia
quotient from \(\xi\). Elsewhere inertia factors through
the residual character quotient. The given finite characters
have conductors dividing
\[
 M=\ell\,2^{n_{\rm II}}\mathcal N.
\]

The resulting inertia image is finite and abelian, and
Frobenius acts trivially on it because all these characters
come from rational characters. The Frobenius-semisimplified
local parameter is thus a sum of Weil lines with zero
monodromy, and each smooth unit character is trivial on
\(1+M\Z_p\). At \(\ell\) the graded Weil lines just
computed have smooth unit characters from the finite
inertia factors of the rational flag, also of conductor
dividing \(M\). These conductor and local conditions enter the rank-four packet comparison. \hfill\mbox{\cite{newton-thorne-unitary}}\label{cor:ordinary-count-ii-point-output-4}

\end{proof}

\subsubsection{Descent and admissible quadratic restrictions}
\label{par:count-two-descent}

The rational local calculation fixes the level. To use the count of cuspidal representations, one must still verify descent and cuspidality after each imaginary quadratic restriction allowed in that count.

\begin{proposition}[Quadratic packet restrictions of the even-order target]\label{prop:count-two-quadratic-packets}
After the stated descent to \(K\), every admissible imaginary
quadratic restriction on the fixed packet list has the regular
algebraic cuspidal attachment at the level \(M\) from the rational local calculation. Its ordered weights are 
\[
 (j+1,j,1,0).
\]

\end{proposition}

\begin{proof}
\proofstep{Descend to the fitted quadratic field.}

Let \(K\subset L\) be the imaginary quadratic field of the
fitted tower. We descend the untwisted transfer to \(K\)
along a chain of cyclic prime degree extensions inside the
abelian extension \(L/K\). At each step the upper cuspidal
transfer is invariant under the cyclic Galois group:
on the Galois side conjugation carries the restriction to an isomorphic representation, since
\[
 r(ghg^{-1})=r(g)r(h)r(g)^{-1}.
\]
Thus it preserves the good Frobenius parameters. Strong multiplicity one proves invariance.
Cyclic cuspidal descent then gives a cuspidal representation
over the lower field. Its weight is regular algebraic,
as can be checked after base change at infinity. Apply
unramified Galois attachment for regular algebraic cuspidal
representations over CM fields to this descent. Its
semisimple representation has the same restriction to
the upper field as \(r\), by unramified compatibility and
Chebotarev. This restriction is absolutely irreducible
since the restriction to \(G_L\) is.

After identifying the restrictions, let \(r_1,r_2\)
be the two extensions to the lower group \(G\), and let
\(H\triangleleft G\) be the upper group. For \(g\in G\) put
\[
 C_g=r_1(g)r_2(g)^{-1}.
\]
Write \(r_H\) for the common restriction. Normality of \(H\) gives, for \(h\in H\),
\[
 C_g r_H(h)C_g^{-1}=r_1(g)r_H(g^{-1}hg)r_1(g)^{-1}=r_H(h).
\]
By irreducibility and Schur's lemma this ratio is a
scalar \(\zeta(g)\). Writing \(r_1(g)=\zeta(g)r_2(g)\) gives
\[
 \zeta(gh)r_2(gh)=r_1(g)r_1(h)=\zeta(g)\zeta(h)r_2(gh).
\]
Hence \(\zeta\) is a character, and \(\zeta|_H=1\)
makes it a finite character of the cyclic quotient.
Twist the descended cuspidal representation by the inverse
finite character.
Its attachment is now \(r\) restricted to the lower field.
The polarization of \(r\) and strong multiplicity one give
the corresponding essential conjugate duality on the
descended cuspidal representation. This proves the assertion
at every stage with the same normalization.
\hfill\mbox{\cite{clozel-rajan,harris-lan-taylor-thorne}}\label{cor:ordinary-count-ii-point-output-5}

\proofstep{Pass to each admissible quadratic field.}

The count requires cuspidality and the level bound over each admissible field.


Fix its eighteen imaginary quadratic fields of pairwise coprime discriminants.
In choosing the ramified prime of \(K\) in
Paragraph~\ref{par:count-ii-tower}, avoid these fixed
discriminants in addition to the rational ramification set.
Dirichlet's theorem in that paragraph allows any such finite
exclusion. For \(K'\) on the list with \(r|_{G_{K'}}\) irreducible,
write \(\eta_{KK'/K'}\) for the quadratic character of
\(KK'/K'\). We use
\[
        r|_{G_{K'}}\not\simeq
               r|_{G_{K'}}\otimes\eta_{KK'/K'} .
\]
This suffices for irreducibility on the index-two subgroup. Indeed, reducibility on that index-two group
would give the quadratic self-twist of \(G_{K'}\)
corresponding to \(KK'/K'\).

At the chosen ramified prime
of \(K\), inertia lies in \(G_{K'}\), acts trivially on \(r\)
and takes the value \(-1\) on the quadratic character.
For such an inertia element \(s\), if \(\nu\) is the quadratic character, then
\[
 \operatorname{tr}r(s)=\dim r,\qquad
 \operatorname{tr}((r\otimes\nu)(s))=-\dim r .
\]
These characteristic-zero traces distinguish the two representations. Cyclic base change
from \(K\) to \(KK'\) is cuspidal by irreducibility there,
and the same descent to \(K'\) now gives the required
attachment. The argument applies to any specified finite
list by the same finite exclusion. No degree or discriminant
of the intermediate tower enters the packet bound.

At a place \(w\mid p\) of \(K'\), restricting the rational
Weil lines gives the inertia characters on the descended
side. Their smooth unit characters factor through the
local norm on units. If \(x\in1+M\mathcal O_{K'_w}\), multiplication by \(x\)
on \(\mathcal O_{K'_w}\) has matrix \(I+MA\) over \(\Z_p\).
Thus
\[
           N_{K'_w/\Q_p}(x)=\det(I+MA)\in1+M\Z_p.
\] Away from \(\ell\)
local--global compatibility for the polarized cuspidal
representation identifies the full Frobenius-semisimplified
Weil--Deligne parameters. Above \(\ell\), compatibility up
to semisimplification suffices. All ratios of the Weil
lines have absolute value one on Frobenius, so the
semisimplification admits only zero monodromy. Hence the
local parameters on the descended side are sums of Weil
lines and their inertia characters have the asserted bounds.
Under local Langlands the local representation is a
principal series subquotient.

It is generic by cuspidality,
and a generic subquotient of a reducible induction of
linked rank-one characters has a segment of length greater
than one and nonzero monodromy. So it is the irreducible
principal series of those lines. Let \(K(M)\) be the principal congruence subgroup. If \(\chi_B\) denotes the full inducing character of the Borel \(B\), including its modulus factor, then \(\chi_B|_{B\cap K(M)}=1\). Hence
\[
 f(bk)=\chi_B(b)\quad(b\in B,\ k\in K(M)),\qquad
 f|_{G\setminus BK(M)}=0
\]
is well-defined and right \(K(M)\)-invariant, with \(f(1)=1\). Where
\(p\nmid M\) the inducing characters are unramified and
there is a spherical vector. The weights after descent
are \(j+1,j,1,0\), from the local de Rham flag. This gives
the quadratic restriction at level \(M\) used in the
packet count.
\hfill\mbox{\cite{blggt,blggt-lg,zelevinsky}}\label{cor:ordinary-count-ii-point-output-6}
\end{proof}
